% Source: https://artofproblemsolving.com/community/c3h2505412_what_does_this_usually_mean
% Stable identity: aops-topic-2505412-q01
\textbf{(a)} By unique factorisation, the divisors of $n$ are exactly the numbers $p_1^{f_1}\cdots p_k^{f_k}$ with $0\le f_i\le e_i$, and different exponent choices give different divisors. There are $e_i+1$ choices for each $f_i$, so $\tau(n)=(e_1+1)\cdots(e_k+1)$.

\textbf{(b)} $2021=43\cdot47$, so $\tau(2021)=2\cdot2=\boxed{4}$. Counting prime factors in $10!=1\cdot2\cdots10$ gives $10!=2^8\cdot3^4\cdot5^2\cdot7$, so $\tau(10!)=9\cdot5\cdot3\cdot2=\boxed{270}$.

\textbf{(c)} We need $\prod(e_i+1)=12$, so the exponent pattern (in decreasing order) is $(11)$, $(5,1)$, $(3,2)$ or $(2,1,1)$. The smallest number with each pattern puts the largest exponents on the smallest primes: $2^{11}=2048$, $2^5\cdot3=96$, $2^3\cdot3^2=72$, $2^2\cdot3\cdot5=60$. The smallest is $\boxed{60}$.
